\documentclass[11pt,a4paper]{article}
\usepackage[margin=26mm]{geometry}
\usepackage[T1]{fontenc}
\usepackage{lmodern}
\usepackage{amsmath,amssymb,amsthm,mathtools}
\usepackage{microtype}
\usepackage{xcolor}
\usepackage[colorlinks=true,linkcolor=blue!45!black,urlcolor=blue!45!black]{hyperref}
\usepackage{fancyhdr}
\pagestyle{fancy}
\fancyhf{}
\fancyhead[L]{Erd\H{o}s 708: joint estimate addendum}
\fancyhead[R]{23 November 2026}
\fancyfoot[C]{\thepage}
\setlength{\headheight}{14pt}
\newtheorem{theorem}{Theorem}[section]
\newtheorem{lemma}[theorem]{Lemma}
\newtheorem{proposition}[theorem]{Proposition}
\newtheorem{corollary}[theorem]{Corollary}
\newtheorem{claim}[theorem]{Unproved sufficient statement}
\newcommand{\pos}[1]{(#1)_+}
\newcommand{\Ph}{\Phi}
\newcommand{\Om}{\Omega}
\newcommand{\ind}{\mathbf{1}}
\newcommand{\raw}{L_{\mathrm{raw}}}
\title{\vspace{-13mm}Erd\H{o}s 708:\\Exact coverage and a canonical joint certificate}
\author{Research addendum and proof audit}
\date{23 November 2026}
\begin{document}
\maketitle
\vspace{-5mm}
\noindent\textbf{Status.}
This addendum does \emph{not} prove \(J_{\rm weak}\), the unrestricted
asymptotic \(g(n)\le(2+o(1))n\), or the exact \(2n\) bound.
It proves an exact global coverage identity, records a precise conditional
route to the desired contradiction, and proves a canonical signed
certificate whose positive contributions never reuse a target budget.
An explicit counterexample rules out one attempted local completion.
All statements labelled unproved remain unproved. Internal mathematical
audits are not external peer review or formal verification.

\section{Fixed ambient problem and the requested estimate}
Let \(A\subseteq[2,M]\) consist of \(m\) distinct integers, and let
\(I=\{x+1,\ldots,x+M\}\), where \(x\ge0\). Put
\[
 D_p(A)=\sum_{a\in A}v_p(a),\qquad
 X_w(n)=\sum_p w_pv_p(n).
\]
All target valuations in this note are the actual, untruncated valuations.
The boxed fractional covering problem and its dual are
\begin{align}
 \raw(A,I)&=\min\left\{\sum_{b\in I}z_b:
 0\le z_b\le1,\ \sum_{b\in I}z_bv_p(b)\ge D_p(A)\right\},\\
 \Ph_A(w)&=\sum_{a\in A}X_w(a)
       -\sum_{b\in I}\pos{X_w(b)-1}.
\end{align}
An optimal dual can be chosen with \(0\le w_p\le1\), supported on the
source primes. We write \(F=\raw(A,I)=\Ph_A(w)\) for a paired optimum.
Complementary slackness gives
\begin{align}
 X_w(b)>1&\ \Longrightarrow\ z_b=1, &
 X_w(b)<1&\ \Longrightarrow\ z_b=0, \label{eq:cs1}\\
 w_p>0&\ \Longrightarrow\
 \sum_{b\in I}z_bv_p(b)=D_p(A). \label{eq:cs2}
\end{align}

In the previously normalized integer-minimal obstruction at budget
\(c=2+3\delta\), the established reduction supplies
\[
 F>(2+3\delta)m-r,\qquad r=o_\delta(m),
 \tag{N}
\]
where \(r\) is the number of source-prime rows. The source normalization,
the core/external-prime restrictions, and the rounding estimate leading
to (N) belong to the preceding report. They are not strengthened by
assertion here.

In its original notation, the requested joint estimate is
\[
 \sum_{a\in G}\bigl(X_w(a)-2-\delta\bigr)
 \le \sum_{b\in I}\pos{X_w(b)-1}
       +C_\delta\Xi+\delta|G|+o_\delta(m).
 \tag{\(J_{\rm weak}\)}
\]
The retained source set \(G\) and row-excess error \(\Xi=o_\delta(m)\)
come from that normalization. The results below do not establish this
inequality.

\section{An exact global coverage identity}
\begin{theorem}\label{thm:coverage}
For every coordinatewise reduction \(0\le u\le w\),
\begin{align}
 \Ph_A(u)&=\sum_{b\in I}z_b\min\{X_u(b),1\}, \label{eq:coverage}\\
 F-\Ph_A(u)&=\sum_{b\in I}z_b\pos{1-X_u(b)}. \label{eq:loss}
\end{align}
Both identities hold for every \(M,A,x\) and every paired optimum.
There is no restriction on the number or sizes of active primes.
\end{theorem}
\begin{proof}
Since \(u_p=0\) whenever \(w_p=0\), the exact active-row equations
\eqref{eq:cs2} imply
\[
 \sum_{a\in A}X_u(a)=\sum_{b\in I}z_bX_u(b).
\]
If \(X_u(b)>1\), then \(X_w(b)>1\), so \(z_b=1\). Consequently
\[
 \sum_{b\in I}\pos{X_u(b)-1}
 =\sum_{b\in I}z_b\pos{X_u(b)-1}.
\]
Subtracting proves \eqref{eq:coverage}. Also \(z_b>0\) implies
\(X_w(b)\ge1\), so the same formula at \(u=w\) yields
\(F=\sum_bz_b\). Subtraction proves \eqref{eq:loss}.
\end{proof}

The loss at a target is at most its single chosen mass \(z_b\).
Thus overlapping prime blocks cannot spend the same target repeatedly
when the loss is estimated through \eqref{eq:loss}.
For example, exact row conservation gives, for any block \(E\) of active primes,
\[
 \sum_b z_b\Om_E(b)=\sum_a\Om_E(a).
\]
If the right side is at most \(Km\), targets with \(\Om_E(b)>L\)
contribute at most \(Km/L\) to \eqref{eq:loss}. This assertion concerns
the bounded loss, not the generally unbounded target hinge.

\subsection{A precise sufficient estimate for completion}
Fix \(\delta>0\), put
\[
 \lambda=\frac{2+\delta}{\log M},\qquad
 u_p=\min\{w_p,\lambda\log p\}.
 \label{eq:logcap}\tag{C}
\]
For every source,
\[
 X_u(a)\le\lambda\log a\le2+\delta.
\]
Hence \(\Ph_A(u)\le(2+\delta)m\). Define the single loss functional
\[
 \mathcal E_\delta(A,I;w,z)
   =\sum_{b\in I}z_b
       \left(1-\sum_p
          \min\{w_p,\lambda\log p\}v_p(b)\right)_+.
 \tag{L}
\]
Theorem \ref{thm:coverage} gives the unconditional inequality
\[
 \boxed{F\le(2+\delta)m+\mathcal E_\delta(A,I;w,z).}
 \tag{R}
\]

\begin{claim}\label{claim:cap}
On the entire actual integer-minimal class used in (N), uniformly in
\(M,A,x,w,z\),
\[
 \mathcal E_\delta(A,I;w,z)\le\delta m+o_\delta(m).
 \tag{U}
\]
\end{claim}
If (U) were proved, (R) would give
\[
 (2+3\delta)m-o_\delta(m)<F
 \le(2+2\delta)m+o_\delta(m),
\]
which is impossible for fixed \(\delta>0\) and large \(m\).
This is an exact conditional contradiction. Neither the coverage
identity nor the bound \(X_u(a)\le2+\delta\) proves (U).
The issue is precisely the chosen mass of targets whose capped load
falls below one.

There is a strictly weaker sufficient quantifier: for each actual
minimal instance it is enough that \emph{one} optimal dual vector \(w\)
satisfy (U), with an error uniform in the instance. For fixed \(w\),
the loss in (L) is independent of the choice of paired optimal \(z\),
because it equals \(F-\Ph_A(u)\). This existential version is also
unproved; it must not be inferred from the proved identities.

\subsection{Why a uniform statement without minimality is false}
Let \(p\) be fixed, let \(m=\lfloor M/p\rfloor\), and take
\(A=\{p,2p,\ldots,mp\}\), \(I=[1,M]\).
The primal vector \(\ind_A\) is feasible with cost \(m\).
The vector \(w_p=1\), all other weights zero, has dual value \(m\),
so this is a paired optimum. For (C), with \(p,\delta\) fixed and
\(M\) sufficiently large,
\[
 \sum_{a\in A}X_u(a)
 =\lambda\log p\sum_{a\in A}v_p(a)
 =O_{p,\delta}(m/\log M).
\]
The identity \eqref{eq:loss} therefore gives
\(\mathcal E_\delta=m-O_{p,\delta}(m/\log M)\).
This is not a counterexample to (U): its sources do not satisfy the
high-load minimal-obstruction conditions. It proves that those
conditions must be used in any valid argument for (U).

\section{What the prime-conserving flow proves exactly}
Assume, in this section only, that every proper source subset satisfies
\(\raw(S,I)\le c|S|\), at the same \(M,I\).
Let \(h_b=\pos{X_w(b)-1}\) and \(H=\sum_bh_b\).
Build a network
\[
 \text{start}\ \longrightarrow\ a\ \longrightarrow\ p\
 \longrightarrow\ b\ \longrightarrow\ \text{end}
\]
with capacities, in that order,
\[
 q_a=\pos{X_w(a)-c},\quad
 w_pv_p(a),\quad w_pv_p(b),\quad h_b.
\]
\begin{proposition}\label{prop:flow}
Its maximum flow is exactly
\[
 \min\left\{\sum_aq_a,H\right\}.
\]
If \(F>cm\), its total deficiency is exactly \(F-cm\).
\end{proposition}
\begin{proof}
For a cut with source subset \(S\) and prime subset \(P\) on the
start side, first optimize the sides of the target vertices.
Sources of quota zero may be removed from \(S\).
The cut deficiency relative to \(\sum_aq_a\) becomes
\[
 \sum_{a\in S}X_{w_P}(a)-c|S|
 -\sum_b\min\{h_b,X_{w_P}(b)\}.
\]
Because \(0\le X_{w_P}\le X_w\), the last sum is at least
\(\sum_b\pos{X_{w_P}(b)-1}\).
The deficiency is therefore at most
\(\Ph_S(w_P)-c|S|\). It is nonpositive for proper \(S\), and is
at most \(F-cm\) for \(S=A\).
If \(F>cm\), deletion of one source gives
\(X_w(a)\ge c+(F-cm)>c\), so
\(\sum_aq_a=H+(F-cm)\).
The full-source cut has capacity \(H\). Max-flow/min-cut now proves
both assertions. If \(F\le cm\), every quota is routable.
\end{proof}

Thus this flow representation adds no unaccounted linear gain.
Its unknown global deficiency is the very quantity that needs to be
excluded. A path preserving one prime is not automatically a path
through a hot common divisor, and it does not by itself prove a
two-target cover.

\section{A canonical signed certificate with one positive root}
This section does not require optimality. Fix a finite prime set \(P\)
and arbitrary weights \(0<w_p\le1\). Order the prime occurrences of
an integer increasingly, with multiplicity.
For an integer of load greater than one, let \(B\) be its first
prefix crossing one. If \(p_0\) is the largest prime of \(B\), then
\[
 X(B/p_0)\le1<X(B),\qquad h_B=X(B)-1\in(0,1].
 \tag{B}
\]
The remaining ordered pieces after deleting the first total unit are
\(h_B\) and the weights of all subsequent prime occurrences.
For a current piece \(u\in(0,1]\) and right-tail load \(R\ge0\), put
\[
 \psi_u(R)=\pos{u+R-1}-\pos{R-1}
          =\min\{u,\pos{u+R-1}\}\le uR.
 \tag{P}
\]
Telescoping gives exactly
\[
 \pos{X(n)-2}=\sum_{\text{pieces }u}\psi_u(R_u).
 \tag{T}
\]

For each root \(B\), use a base carrier \(d=B,u=h_B,q=p_0\);
also use carriers \(d=Bq^j,u=w_q\), where \(q\in P\) and \(q\ge p_0\)
and \(j\ge1\).
Keep only carriers with \(dq\le M\), since all others have zero
source contribution. Set \(Q_B=\{p\in P:p<p_0\}\) and
\[
 K_{B,d}(b)=\ind_{d\mid b}
             -\sum_{p\in Q_B}\ind_{dp\mid b}.
 \tag{K}
\]
\begin{theorem}\label{thm:star}
For every positive target integer \(b\),
\[
 \sum_{\text{eligible carriers}}uK_{B,d}(b)
 \le\pos{X(b)-1}.
 \tag{S}
\]
There is at most one root with positive contributions at any target.
\end{theorem}
\begin{proof}
A positive kernel equals one. Its conditions say that \(d\mid b\)
and no prime below \(p_0\) divides \(b/d\).
Since \(d/B\) uses only a prime at least \(p_0\), \(B\) must contain
the complete lower-prime blocks of \(b\).
Condition (B) then identifies \(B\) as the unique actual first
crossing root of \(b\). All other roots have nonpositive contributions.
For this unique root, the positive base and star terms sum to at most
\[
 h_B+\sum_{q\ge p_0}w_qv_q(b/B)=X(b)-1.
\]
Eligibility may only remove nonnegative terms. This proves (S),
including targets with arbitrarily large prime powers.
\end{proof}

For one carrier let \(T_d=M/d\). Its exact total source contribution
in (T), summed over the full prefix, is
\[
 S_{B,d}=
 \sum_{\substack{t\le T_d\\
          \gcd(t,\prod_{p\in Q_B}p)=1}}
       \psi_u\!\left(\sum_{\substack{r\in P\\r\ge q}}
                       w_rv_r(t)\right).
 \tag{F}
\]
This formula handles repeated occurrences: for a \(q^j\) carrier,
occurrences below \(q\) in \(t\) precede the current occurrence and
are excluded from its right-tail load. By (P),
\[
 S_{B,d}\le
 uT_d\sum_{\substack{r\in P\\q\le r\le T_d}}\frac{w_r}{r-1}.
 \tag{F1}
\]
The upper cutoff \(r\le T_d\) is valid and can materially improve
the bound. On every shifted interval,
\[
 \sum_{b\in I}K_{B,d}(b)
 \ge T_d\left(1-\sum_{p\in Q_B}\frac1p\right)-|Q_B|-1.
 \tag{F2}
\]
Together these prove full hinge domination if every contributing
carrier satisfies
\[
 \sum_{p\in Q_B}\frac1p+
 \sum_{\substack{r\in P\\q\le r\le T_d}}\frac{w_r}{r-1}
 +\frac{|Q_B|+1}{T_d}\le1.
 \tag{FC}
\]
Condition (FC) is sufficient, not known on the entire minimal class.
In particular its first sum is unweighted. Weighted mean control
alone does not imply it.

\section{An exact failure of a separate-root completion}
Replace (K) by the exact nonnegative kernel
\[
 K^{\rm ex}_{B,d}(b)
 =\ind_{d\mid b}
   \ind_{\gcd(b/d,\prod_{p\in Q_B}p)=1}.
 \tag{E}
\]
The same one-positive-root argument still gives
\(\sum uK^{\rm ex}_{B,d}\le\pos{X-1}\).
However the tempting \emph{separate} comparison
\[
 S_{B,d}\ \stackrel{?}{\le}\
 u\sum_{b\in I}K^{\rm ex}_{B,d}(b)
 \tag{false}
\]
is false, even with weighted reciprocal mean less than one.

Take
\[
 P=\{2,3,5,7,11,13\},\qquad
 w_2=\frac1{10},\quad w_{13}=1,\quad
 w_3=w_5=w_7=w_{11}=\frac1{1000}.
\]
Then \(\sum_pw_p/(p-1)=0.18435\).
For \(B=d=26\), \(p_0=13\), \(u=h_B=1/10\), let
\[
 M=338,\qquad I=\{2939,\ldots,3276\}.
\]
The source cofactor length is \(13\), and the only prefix cofactors
coprime to \(2\cdot3\cdot5\cdot7\cdot11\) are \(1,13\).
Their total in (F) is \(\psi_{1/10}(0)+\psi_{1/10}(1)=1/10\).
The contributing source is \(338=2\cdot13^2\), of load \(2.1\).

The multiples of \(26\) in \(I\), after division by \(26\), are
\(\{114,\ldots,126\}\). Each is divisible by \(2,3,5,7\), or \(11\).
In particular the two less immediate odd cases are
\(119=7\cdot17\) and \(121=11^2\).
Thus the right side of (false) is zero.

This is a counterexample to the individual-fibre comparison, not
to \(J_{\rm weak}\), (U), or the desired asymptotic. The weights
were not claimed to be an optimum of an actual minimal obstruction.
On the targets, additional earlier prime factors move contributions
to other roots. A complete proof must account for that transfer
jointly, with its signs and its single shared budget.

\section{A uniform finite theorem for all low-prime states}
The following auxiliary theorem controls all state combinations at
once, rather than multiplying an error by their number.
\begin{theorem}\label{thm:tv}
Let \(P\) be any \(k\)-element prime set and \(s\ge1\) an integer.
Define \(z_p(n)=\min\{v_p(n),s\}\), \(z(n)=(z_p(n))_{p\in P}\),
and \(\mu=\sum_{p\in P}1/p\).
Let \(N_0(\sigma)\) and \(N_I(\sigma)\) count each state in
\([1,M]\) and \(I\), respectively. Then, uniformly in the shift,
\[
 \sum_\sigma\pos{N_0(\sigma)-N_I(\sigma)}
 \le C_s\bigl(k^s+M\mu^{s+1}+k^{s+1}\bigr).
 \tag{TV}
\]
\end{theorem}
\begin{proof}
For any \(f:\{0,\ldots,s\}^P\to[0,1]\), apply presence-set
M\"obius inversion to \(f\), and retain its terms supported on at
most \(s\) primes. Call the resulting function \(T_sf\).
It agrees with \(f\) when \(\nu(n)=|\{p\in P:p\mid n\}|\le s\).
In general
\[
 |f(z(n))-T_sf(z(n))|
 \le B_s\binom{\nu(n)}{s+1}.
 \tag{TV1}
\]
Indeed the support-\(j\) coefficients have absolute value at most
\(2^j\), and for \(\nu\ge s+1\),
\(\binom{\nu}{j}\le\binom{s+1}{j}\binom{\nu}{s+1}\).

Each support-\(j\) term has a finite prime-power difference expansion
in indicators of moduli \(\prod_{p\in S}p^{e_p}\), \(1\le e_p\le s\).
Its total coefficient absolute sum is at most \((4s)^j\).
All nonconstant retained coefficients therefore have total absolute
sum at most \(A_sk^s\). Every multiple count differs by at most one
between the prefix and an equal-length interval, so
\[
 \left|\sum_{n\le M}T_sf(z(n))-\sum_{b\in I}T_sf(z(b))\right|
 \le A_sk^s.
 \tag{TV2}
\]
The tails on both actual sets satisfy
\begin{align*}
 \sum_{n\le M}\binom{\nu(n)}{s+1}
   &\le M e_{s+1}((1/p)_{p\in P}),\\
 \sum_{b\in I}\binom{\nu(b)}{s+1}
   &\le M e_{s+1}((1/p)_{p\in P})+\binom{k}{s+1}.
\end{align*}
Combine these with (TV1)--(TV2) and
\(e_{s+1}\le\mu^{s+1}/(s+1)!\).
The bound is uniform over \(f\). Choosing the indicator of all
states with \(N_0(\sigma)>N_I(\sigma)\) proves (TV).
\end{proof}

State bins give an injective state-preserving matching of every
source subset except at most the deficit in (TV). This is one
matching through disjoint bins. It does not automatically cover
the remaining prime demands with one additional target per source.

For the low substantial primes in the established sparse reduction,
\[
 s=\lfloor1/\eta\rfloor+1,\quad
 y=M^{1/(s+1)},\quad
 P=\{p\le y:w_p\ge\eta\},\quad \rho=m/M\to0,
\]
the companion absorption proof yields
\(\mu=O_{c,\eta}(\rho^{1/s})\), \(k\le y\mu\).
Consequently (TV) is
\[
 O_{c,\eta}\!\left(m\bigl[\rho^{1/s}+M^{-1/(s+1)}\bigr]\right)
 =o(m).
\]
The companion note also controls the full weighted demand of the
unmatched sources, not only their cardinality. This auxiliary
matching is not presented as a solution of the unrestricted problem.

\section{Remaining obligation}
The global identities above are exact, and their proofs do not
assume a small core, a bounded number of substantial weights, or
gaps between primes. The finite state theorem likewise sums all
states jointly.

The missing step is still a quantitative joint estimate on actual
minimal obstructions. One sufficient form is (U). In the canonical
certificate it is the signed compensation between different roots.
Neither separate-root averaging, a sum of independently chosen
covers, nor the prime-conserving flow alone supplies that estimate.
No counterexample to \(J_{\rm weak}\) has been obtained.

\end{document}
